% !TeX root = ../main.tex
\chapter{Kết luận}
\label{chap:conclusion}

Nhóm đã xây dựng đầy đủ một quy trình Decision Tree cho Breast Cancer Wisconsin
(Diagnostic): kiểm tra dữ liệu, chia train/test có stratification, huấn luyện
baseline, trình bày cây, đánh giá nhiều metric, phân tích rule và thử ba cách
regularization độc lập.

Baseline đạt 91,23\% test accuracy nhưng có training accuracy 100\%, depth 7
và 19 leaf. Khoảng cách train--test cho thấy cây không ràng buộc có nguy cơ
overfit. Giới hạn \texttt{max\_depth=4} và cost-complexity pruning đều tăng
accuracy lên 93,86\%, giảm error rate xuống 6,14\%. Điều chỉnh
\texttt{min\_samples\_leaf=3} làm cây nhỏ hơn nhưng không cải thiện accuracy.

Max Depth được chọn vì giữ malignant recall 92,86\% và có weighted F1 cao nhất
trong hai mô hình đồng hạng accuracy. Pruning vẫn là lựa chọn tốt nếu mục tiêu
ưu tiên cấu trúc nhỏ và dễ giải thích. Kết quả chứng minh regularization có thể
giảm overfitting, nhưng việc chọn mô hình phải dựa trên nhiều metric và chi phí
lỗi thay vì chỉ nhìn accuracy.

\section{Hạn chế}

\begin{itemize}
  \item Dataset có 569 mẫu; chênh lệch một vài mẫu test có thể làm metric thay
  đổi đáng kể.
  \item Thí nghiệm dùng một train/test split cố định để tái lập. Nested
  cross-validation hoặc repeated stratified cross-validation sẽ cho ước lượng
  ổn định hơn về độ biến thiên.
  \item Feature importance của một cây có thể thiên về feature có nhiều khả
  năng chia; permutation importance hoặc SHAP có thể bổ sung góc nhìn.
  \item Đây là thí nghiệm học thuật, không phải hệ thống hỗ trợ chẩn đoán lâm
  sàng. Các threshold học được không nên diễn giải như quy tắc y khoa.
\end{itemize}

\section{Hướng phát triển}

Có thể mở rộng bằng nested cross-validation, đánh giá ROC--AUC, tối ưu trực
tiếp recall malignant, kiểm tra class weighting, feature selection và so sánh
với Random Forest hoặc Gradient Boosting. Với mô hình cây đơn, việc xác nhận
độ ổn định của rule qua nhiều split cũng là hướng phân tích quan trọng.
